A slash inside a word, a path or a URL never opens or closes \emph{italics}.

Bare URLs are autolinked whole: see \href{http://a/b/c/}{http://a/b/c/} for this, and
\href{https://example.com/a/b/?q=1\&r=/x/\#frag/}{https://example.com/a/b/?q=1\&r=/x/\#frag/} with a query and a fragment, and
\href{ftp://files.example.org/pub/x/}{ftp://files.example.org/pub/x/} too.

A URL ending a sentence keeps the full stop outside: read \href{http://a.com/b/c/}{http://a.com/b/c/}.

Already fine before: \href{http://www.example.com/a/b/}{www.example.com/a/b/} and \href{http://a/b/c/}{http://a/b/c/}.

Paths stay paths: see /usr/local/bin/ here, (/usr/local/bin/) in parentheses,
~/notes/ and ./src/ and ../lib/ and C:/Users/x/ on other systems.

Words and numbers: either and/or both, and the fraction 1/2/3.

A URL and real italics together: see \href{http://a.com/}{http://a.com/} and \emph{this phrase} here.

